\chapter{Ovarian Cystectomy}

\section*{Introduction \& Clinical Indications}
Ovarian cystectomy removes a cyst while preserving ovarian tissue when this is safe and appropriate. Many simple functional cysts resolve without surgery. Surgery may be considered for persistent or enlarging cysts, significant symptoms, torsion, endometrioma or dermoid cyst in selected patients, or a mass whose features require tissue diagnosis. A mass suspicious for malignancy should be managed with appropriate gynecologic-oncology input; cystectomy is not automatically the correct operation.

The choice between observation, cystectomy, oophorectomy, or staging surgery depends on symptoms, ultrasound features, size and persistence, age and menopausal status, fertility goals, family history, and malignancy risk. Pregnancy requires a separate individualized assessment; most simple cysts are observed unless complications or malignancy are suspected.

\section*{Relevant Surgical Anatomy}
The ovary lies beside the uterus and is connected by the ovarian ligament and suspensory ligament. The ovarian vessels run in the suspensory ligament and communicate with uterine vessels. The ureter lies retroperitoneally near the pelvic sidewall and must be identified when dissection is extensive. Preservation of ovarian cortex, vascular supply, and the contralateral ovary is important when fertility or endocrine function matters.

\section*{Preoperative Workup \& Preparation}
Evaluation begins with history, examination, pregnancy testing when relevant, and high-quality transvaginal or transabdominal ultrasound. MRI can further characterize an indeterminate mass. Tumor markers are selected according to age, imaging, and clinical context; CA-125 is not a stand-alone screening or diagnostic test and may be elevated in benign conditions. A mass with concerning features should be referred or discussed with a gynecologic oncologist before surgery.

Consent includes observation and alternative operations, possible conversion to laparotomy or oophorectomy, rupture or spillage, bleeding, injury to bowel, bladder, ureter, or vessels, reduced ovarian reserve, adhesions, recurrence, and the possibility of unexpected malignancy. Routine antibiotic and laboratory decisions follow the operation and patient risk. Pregnancy, uncontrolled coagulopathy, and serious systemic illness require individualized planning rather than being treated as universal absolute contraindications.

\section*{Surgical Technique \& Procedure}
\begin{enumerate}
  \item \textbf{Access}: Under general anesthesia, the patient is positioned for laparoscopy or laparotomy. The abdomen and pelvis are inspected for the cyst, torsion, endometriosis, adhesions, or signs of malignancy.
  \item \textbf{Exposure}: The ovary and nearby ureter and vessels are identified. The approach is chosen to avoid unnecessary rupture, particularly when malignancy is possible.
  \item \textbf{Cyst removal}: The cyst is opened or dissected from the ovarian cortex using tissue-sparing technique. Hemostasis is achieved with the least thermal injury practical; the cyst is removed in a protected manner when appropriate.
  \item \textbf{Completion}: The specimen is sent for histopathologic examination. The pelvis is irrigated if needed, hemostasis is confirmed, and the incisions are closed. If cancer is suspected, the planned operation follows oncologic principles rather than simple cystectomy.
\end{enumerate}

\section*{Complications \& Risks}
Risks include bleeding, infection, anesthesia complications, venous thromboembolism, injury to bowel, bladder, ureter, or vessels, conversion to open surgery, cyst rupture and chemical peritonitis, adhesions, recurrent cyst, and reduced ovarian reserve or ovarian failure. Endometrioma surgery can reduce ovarian reserve, so fertility goals and alternatives should be discussed.

\section*{Postoperative Care \& Recovery}
Patients receive pain control, wound instructions, early mobilization, and guidance on activity and sexual intercourse. The final pathology result determines whether additional evaluation is needed. Urgent review is required for fever, worsening pain, heavy bleeding, persistent vomiting, urinary difficulty, leg swelling, or breathing symptoms. Follow-up is individualized and may include review of pathology, symptoms, fertility goals, and repeat imaging when clinically indicated.

\section*{Outcomes \& Prognosis}
Recovery is often faster after minimally invasive surgery for a benign-appearing mass, but outcomes depend on the cyst type, operative findings, ovarian reserve, and underlying disease. Symptom recurrence is possible, particularly with endometriosis or functional cysts. No single success rate applies to all ovarian cystectomies.

\section*{Additional Clinical Considerations}
The preoperative question is not simply whether a cyst is large. Ultrasound morphology, vascularity, septations, papillary projections, solid components, ascites, bilaterality, and growth pattern influence the risk assessment. In a premenopausal patient, a simple cyst may be observed with repeat imaging, while a complex or persistent mass requires a more deliberate plan. In a postmenopausal patient, the threshold for specialist assessment is lower because the differential diagnosis and malignancy risk differ. Torsion is a time-sensitive clinical diagnosis; imaging can support it but a normal Doppler study does not reliably exclude intermittent torsion.

When malignancy is plausible, the operation should be planned with a gynecologic oncologist. Avoiding rupture and spillage is desirable, but an attempt to preserve the ovary must not compromise staging or oncologic safety. The specimen should be retrieved in a protected bag when feasible, and the surgical team should have a plan for frozen section or conversion if the findings are unexpected. Cyst fluid alone is not a substitute for complete histopathologic assessment. If a malignant or borderline tumor is identified, further surgery, staging, genetic counseling, or systemic treatment may be required.

Ovarian preservation has reproductive and endocrine importance, but cystectomy itself can reduce ovarian reserve, especially when an endometrioma is bilateral, recurrent, large, or treated with extensive coagulation. The surgeon should use a tissue-sparing dissection and minimize thermal injury while achieving hemostasis. Patients pursuing fertility may benefit from a reproductive endocrinology discussion before surgery, particularly when the contralateral ovary is abnormal, the ovarian reserve is already low, or an oophorectomy is possible.

Endometriosis requires counseling that removing an endometrioma does not eradicate endometriosis elsewhere in the pelvis. Pain may persist or recur, and postoperative hormonal suppression may be considered when pregnancy is not immediately desired. Dermoid cysts can contain mature tissues and may cause chemical peritonitis if ruptured; meticulous containment and irrigation are used when spillage occurs. Functional cysts can recur because the ovary remains hormonally active, so a successful operation is not the same as permanent prevention of every future cyst.

After surgery, pathology review is essential. The report may distinguish a functional cyst, endometrioma, mature cystic teratoma, borderline tumor, or invasive malignancy and may determine the need for additional imaging or referral. Patients should receive clear instructions about the expected recovery and urgent symptoms. Severe unilateral pain, fever, fainting, heavy bleeding, persistent vomiting, urinary retention, calf swelling, or shortness of breath may indicate a complication requiring prompt evaluation. Follow-up also addresses menstrual function, fertility plans, ovarian reserve when relevant, and whether repeat imaging is needed.

\section*{Additional Follow-Up Considerations}
The final diagnosis determines whether the patient returns to routine gynecologic care or needs specialist follow-up. A benign functional cyst may need no further imaging, whereas endometriosis, a dermoid, a borderline tumor, or an indeterminate mass may require surveillance. The pathology report should be explained in plain language, including whether the cyst was completely removed, whether the ovary was preserved, and whether further treatment or genetic assessment is recommended.

Fertility counseling is individualized. Ovarian reserve can be assessed with antral follicle count, anti-Mullerian hormone, or other tests when the result will change planning, but no single marker predicts natural conception by itself. Patients who wish to conceive should be told when intercourse or assisted reproduction can resume and should seek care for persistent pain, irregular bleeding, or difficulty conceiving. Endometriosis and recurrent cysts may require medical suppression or a fertility-preservation discussion.

\section*{Fertility and Pathology Follow-Up}
Preserving the ovary does not necessarily preserve the same ovarian reserve as before surgery. The effect is influenced by cyst type, bilateral disease, repeat operations, hemostatic method, and the amount of cortex removed. A patient with a low reserve, a single functioning ovary, or a strong fertility goal may benefit from counseling about egg or embryo preservation before an elective operation when time permits. The decision is balanced against torsion, malignancy, pain, and other indications that cannot safely be delayed.

Pathology review should distinguish a benign lesion from a borderline or malignant tumor. Unexpected malignancy may require referral for staging rather than an immediate unplanned second operation. The patient should receive a clear plan for repeat imaging, tumor-marker interpretation, genetic counseling, and surveillance when indicated. A normal postoperative scan does not eliminate the possibility of a new cyst, endometriosis, adhesions, or a separate cause of pelvic pain.

\section*{Long-Term Surveillance}
A new or persistent mass after surgery should be evaluated on its own merits. The prior pathology and operative report help determine whether observation, repeat ultrasound, MRI, or oncology assessment is appropriate. Patients with endometriosis may need coordinated pain, hormonal, fertility, and bowel or bladder care. Those with a borderline or malignant lesion require surveillance and treatment according to stage and specialist guidance rather than routine benign-cyst follow-up.

The operation can affect ovarian reserve, but leaving the ovary in place generally preserves endocrine function better than oophorectomy. The decision to operate on the contralateral ovary is especially important in bilateral disease. Fertility preservation, pregnancy timing, and the possibility of assisted reproduction should be discussed before repeat surgery when clinically feasible.

\section*{Patient Education}
Patients should know the expected recovery, the pathology plan, and which symptoms require urgent review. Sudden severe unilateral pain may represent torsion even after a prior cystectomy. Fever, fainting, heavy bleeding, progressive abdominal distension, persistent vomiting, or shortness of breath should not be managed as routine postoperative discomfort. Fertility and contraception counseling are revisited when the pathology and operative findings are available.

\section*{Documentation and Safety}
The operative note should describe the cyst, ovary, contralateral adnexa, rupture or spillage, hemostatic method, estimated preservation of ovarian tissue, and specimens submitted. Patients should receive the pathology result and a plan for follow-up before the episode of care is closed. A persistent or recurrent mass is evaluated according to its current features rather than assumed to be the original benign cyst.

\section*{Practical Follow-Up}
Recovery is assessed by pain, bleeding, wound condition, return of activity, and review of pathology. Persistent or recurrent symptoms should prompt targeted imaging rather than an assumption that the original cyst has returned. Fertility, ovarian reserve, endometriosis treatment, and contraception are revisited according to the patient's goals.

\section*{Safety Summary}
Observation, cystectomy, oophorectomy, and oncologic surgery are different choices. Imaging, age, symptoms, pregnancy status, fertility goals, and malignancy risk determine which option is safest; unexpected pathology should trigger appropriate specialist follow-up.

\section*{Final Safety Note}
The pathology result closes neither the clinical episode nor the fertility discussion. Follow-up is based on the lesion, ovarian reserve, symptoms, pregnancy plans, and any possibility of borderline or malignant disease.

\section*{Completion Checklist}
Confirm pathology communication, symptom review, fertility or contraception counseling, and the need for imaging or specialist follow-up. Sudden severe pain, fever, fainting, heavy bleeding, or persistent vomiting requires urgent evaluation.

\section*{Handoff}
The patient should receive the pathology plan, fertility or contraception advice, recovery instructions, and a clear route for urgent symptoms.

\section*{References}
\begin{enumerate}
  \item American College of Obstetricians and Gynecologists. Evaluation and Management of Adnexal Masses. Practice Bulletin No. 174. Reaffirmed 2025.
  \item American College of Obstetricians and Gynecologists. Ovarian Cysts. Patient FAQ.
  \item Royal College of Obstetricians and Gynaecologists. Management of Suspected Ovarian Masses in Premenopausal Women. Green-top Guideline No. 62.
\end{enumerate}

\clearpage
